\chapter{實驗設計與結果}\label{chap:results}
\begingroup
\setlength{\emergencystretch}{2em}

\section{統計分析資料與比較範圍}

本章針對實驗一的固定年份集成比較，以及實驗二的 FS 集成比較，呈現既有預測結果與公司群聚配對檢定。兩項比較使用同一份資料切分：Old 擬合期間為 1999--2011 年，New 擬合期間為 2012--2013 年，2014 年分為 DSEL 與 Calibration，2015--2018 年為 Test。Test 含 12,282 筆公司年度紀錄，其中 287 筆屬破產正類，來自 3,700 家公司，其中 119 家為破產公司。

完整資料包含十個種子與九種特徵設定，共 90 組已完成的實驗設定。統計分析重新核對每組設定的測試紀錄識別與標籤，再從保存的預測機率與固定門檻重算 AUC、F1-score 及 Recall；本章比較方法所涉及的 560 個模型輸出均與原紀錄一致。此處不重新訓練、不調整門檻，也不依 Test 挑選 FS 方法或靜態成員。

表中 pooled 指標先於每個種子內合併四個測試年度計算，再跨十個種子彙整；種子標準差描述程序隨機變異。差值的信賴區間與近似檢定則使用第~\ref{sec:method-statistics}~節的公司群聚 bootstrap，將同公司跨年的紀錄一併重抽。兩項實驗共 165 項比較共同採 Holm 校正，與只將十個種子視為獨立樣本的檢定具有不同推論單位。

實驗一的歷史 Baseline 使用年份邊界與年內驗證，並未與本章固定年份集成完成同協定的重新調參比較。故本章統計檢定限於下列固定年份方法，不將其與歷史 Baseline 的分數差作為集成效益之檢定證據。

\section{實驗一：固定年份集成比較}\label{sec:results-experiment1}

本節先呈現未進行特徵選取時的集成結果。Old3 與 New3 分別整合單一資料期間的三種採樣模型，Equal6 整合全部六模型，V-static 則依各種子的 Calibration AP，從 49 個混合子集中選定成員。動態方法均使用相同六模型池與 DSEL，因此下列比較具有相同資料與評估協定。

\begingroup
\small\setstretch{1.12}\setlength{\tabcolsep}{3pt}
\setlength{\LTleft}{\fill}\setlength{\LTright}{\fill}
\begin{longtable}{@{}>{\raggedright\arraybackslash}p{0.22\linewidth}>{\raggedright\arraybackslash}p{0.24\linewidth}>{\raggedright\arraybackslash}p{0.24\linewidth}>{\raggedright\arraybackslash}p{0.24\linewidth}@{}}
\caption{實驗一固定年份集成之 pooled 結果}\label{tab:stats-exp1-descriptive}\\
\toprule
方法 & AUC & F1-score & Recall \\
\midrule\endfirsthead
\multicolumn{4}{c}{表~\thetable（續）}\\
\toprule
方法 & AUC & F1-score & Recall \\
\midrule\endhead
\midrule\multicolumn{4}{r}{續下頁}\\\endfoot
\bottomrule\endlastfoot
Old3 & 0.7134\(\pm\)0.0046 & 0.1168\(\pm\)0.0046 & 0.2383\(\pm\)0.0759 \\
New3 & 0.7959\(\pm\)0.0027 & 0.2007\(\pm\)0.0090 & 0.2268\(\pm\)0.0361 \\
Equal6 & 0.7815\(\pm\)0.0026 & 0.1815\(\pm\)0.0065 & 0.2690\(\pm\)0.0414 \\
V-static & 0.8152\(\pm\)0.0040 & 0.2355\(\pm\)0.0159 & 0.2481\(\pm\)0.0233 \\
DCS-OLA & 0.7386\(\pm\)0.0078 & 0.1542\(\pm\)0.0078 & 0.2909\(\pm\)0.0607 \\
DCS-LCA & 0.7369\(\pm\)0.0087 & 0.1631\(\pm\)0.0083 & 0.3195\(\pm\)0.0354 \\
KNORA-E & 0.7613\(\pm\)0.0052 & 0.1837\(\pm\)0.0058 & 0.2714\(\pm\)0.0345 \\
KNORA-U & 0.7481\(\pm\)0.0079 & 0.1817\(\pm\)0.0081 & 0.2620\(\pm\)0.0428 \\
\end{longtable}
\noindent{\footnotesize 數值為十個種子的平均值與種子標準差；標準差不等同公司抽樣之標準誤。V-static 為 Calibration AP 選定的混合靜態方案；Equal6 為六模型等權平均。}
\endgroup


各方法相對 Equal6 的 AUC 差值、公司群聚信賴區間與 Holm 校正值列於表~\ref{tab:stats-exp1-auc}。正差值代表相對基準的 AUC 增加，負差值代表降低。區間為逐項 95\% 區間，顯著性則依全體比較校正後的 \(p_H\) 判斷。

\begingroup
\small\setstretch{1.12}\setlength{\tabcolsep}{3pt}
\setlength{\LTleft}{\fill}\setlength{\LTright}{\fill}
\begin{longtable}{@{}>{\raggedright\arraybackslash}p{0.25\linewidth}>{\raggedright\arraybackslash}p{0.17\linewidth}>{\raggedright\arraybackslash}p{0.33\linewidth}>{\raggedright\arraybackslash}p{0.17\linewidth}@{}}
\caption{實驗一各集成方法相對 Equal6 之 AUC 配對檢定}\label{tab:stats-exp1-auc}\\
\toprule
方法 & \(\Delta\mathrm{AUC}\) & 95\% 信賴區間 & \(p_H\) \\
\midrule\endfirsthead
\multicolumn{4}{c}{表~\thetable（續）}\\
\toprule
方法 & \(\Delta\mathrm{AUC}\) & 95\% 信賴區間 & \(p_H\) \\
\midrule\endhead
\midrule\multicolumn{4}{r}{續下頁}\\\endfoot
\bottomrule\endlastfoot
Old3 & -0.0681 & [-0.0890, -0.0488] & 0.0033\(^{*}\) \\
New3 & +0.0144 & [-0.0019, +0.0313] & 1.0000 \\
V-static & +0.0337 & [+0.0198, +0.0480] & 0.0063\(^{*}\) \\
DCS-OLA & -0.0430 & [-0.0634, -0.0233] & 0.0186\(^{*}\) \\
DCS-LCA & -0.0446 & [-0.0663, -0.0242] & 0.0186\(^{*}\) \\
KNORA-E & -0.0202 & [-0.0387, -0.0029] & 1.0000 \\
KNORA-U & -0.0334 & [-0.0517, -0.0161] & 0.0626 \\
\end{longtable}
\noindent{\footnotesize 差值為候選方法減比較基準；信賴區間為公司群聚 bootstrap 的逐項 95\% 區間，非同時區間。\(p_H\) 為兩項實驗全部 165 項比較共同採 Holm 校正之雙尾近似 \(p\) 值；星號表示 \(p_H<0.05\)。所有比較使用相同公司重抽權重，模型與驗證門檻固定。}
\endgroup


V-static 的 AUC 與 F1-score 相對 Equal6 分別增加 0.0337 與 0.0540，Holm 校正後 \(p_H\) 分別為 0.00632 與 0.02416。AUC 差值的逐項 95\% 信賴區間為 [0.0198, 0.0480]。其靜態成員由 Calibration 決定，因此該差異屬於驗證選模程序與完整六模型等權平均之比較。相對地，Old3 的 AUC 與 F1-score 差值分別為 \(-0.0681\) 與 \(-0.0647\)，兩者 \(p_H\) 均為 0.00330。

DCS-OLA 與 DCS-LCA 的 AUC 相對 Equal6 分別降低 0.0430 與 0.0446，兩者 \(p_H=0.01860\)。New3、KNORA-E 與 KNORA-U 的 AUC 比較未通過整體 Holm 校正；所有方法相對 Equal6 的 Recall 比較亦未達校正後顯著。這些結果顯示，不能預設動態選擇或跨期間整合必然優於所有控制組。

F1-score 與 Recall 的完整檢定列於附錄~\ref{sec:stats-full}。AUC 為主要報告指標，但其改善不自動代表指定門檻下的 Recall 也同時改善；三項指標須分別解讀。

\section{實驗二：FS 集成比較}\label{sec:results-experiment2}

本節以相同種子、測試公司與方法名稱配對，比較八種 FS 條件與 no-FS 結果。50\% 與 80\% 設定分別保留 9 與 14 個特徵。各 FS 條件均重新建立模型池，並沿用原本的驗證選模與門檻規則；本節不依測試分數選出單一 FS 方案。

\begingroup
\small\setstretch{1.12}\setlength{\tabcolsep}{3pt}
\setlength{\LTleft}{\fill}\setlength{\LTright}{\fill}
\begin{longtable}{@{}>{\raggedright\arraybackslash}p{0.16\linewidth}>{\raggedright\arraybackslash}p{0.123\linewidth}>{\raggedright\arraybackslash}p{0.123\linewidth}>{\raggedright\arraybackslash}p{0.123\linewidth}>{\raggedright\arraybackslash}p{0.123\linewidth}>{\raggedright\arraybackslash}p{0.123\linewidth}>{\raggedright\arraybackslash}p{0.123\linewidth}@{}}
\caption{實驗二各 FS 設定相對同方法 no-FS 之 AUC 差值}\label{tab:stats-exp2-auc}\\
\toprule
FS & Equal6 & V-static & OLA & LCA & KNORA-E & KNORA-U \\
\midrule\endfirsthead
\multicolumn{7}{c}{表~\thetable（續）}\\
\toprule
FS & Equal6 & V-static & OLA & LCA & KNORA-E & KNORA-U \\
\midrule\endhead
\midrule\multicolumn{7}{r}{續下頁}\\\endfoot
\bottomrule\endlastfoot
MI 50\% & -0.0171 & -0.0327 & -0.0257 & -0.0431 & -0.0274 & -0.0220 \\
MI 80\% & +0.0007 & -0.0007 & -0.0066 & -0.0184 & -0.0100 & +0.0076 \\
CART 50\% & -0.0304 & -0.0221 & -0.0299 & -0.0295 & -0.0435\(^{*}\) & -0.0229 \\
CART 80\% & -0.0003 & -0.0013 & +0.0026 & +0.0003 & -0.0056 & +0.0013 \\
SHAP 50\% & -0.0220 & -0.0182 & -0.0033 & -0.0109 & -0.0216 & -0.0141 \\
SHAP 80\% & -0.0000 & -0.0036 & -0.0041 & -0.0090 & -0.0129 & -0.0013 \\
RFE 50\% & -0.0198 & -0.0141 & -0.0278 & -0.0286 & -0.0378 & -0.0125 \\
RFE 80\% & +0.0022 & -0.0002 & +0.0115 & -0.0013 & -0.0014 & +0.0087 \\
\end{longtable}
\noindent{\footnotesize 50\% 與 80\% 設定實際保留 9 與 14 個特徵。正值表示 AUC 增加。星號定義與表~\ref{tab:stats-exp1-auc} 相同；完整區間、F1-score、Recall 與校正值見附錄~\ref{sec:stats-full}。}
\endgroup


表~\ref{tab:stats-exp2-auc} 呈現各 FS 條件相對同方法 no-FS 的 AUC 差值。完整的 48 組 AUC 比較，以及對應的 F1-score、Recall 檢定與信賴區間，均列於附錄~\ref{sec:stats-full}。差值的方向與顯著性分開判讀，未達顯著不代表已證明方法等效。

在本次整體校正下，沒有 FS 條件呈現顯著的 AUC 增加，F1-score 的全部 48 項比較亦未達顯著。CART 50\% 搭配 KNORA-E 的 AUC 差值為 \(-0.0435\)，其 \(p_H=0.03900\)，方向為下降；該值接近顯著門檻且受有限次 Monte Carlo 誤差影響，因此不將此單一門檻判斷視為穩健結論。

Recall 共有九項比較通過 Holm 校正，其中 V-static 在 MI 80\%、SHAP 50\%、SHAP 80\% 及 RFE 80\% 下增加，其餘五項為 DCS 方法下的下降。例如，SHAP 50\% 的 V-static Recall 增加 0.1021，\(p_H=0.00632\)；CART 50\% 的 DCS-LCA Recall 降低 0.1847，\(p_H=0.00330\)。此差異顯示 FS 的影響與集成方法有關，不能僅因部分 Recall 改善，就推論 FS 全面提升區辨能力或精確率與召回率的平衡。

對 V-static 而言，FS 後的靜態成員與門檻可依 Calibration 重新決定，其差異反映整體程序的影響。對固定方法而言，F1-score 與 Recall 仍可能受到 FS 後的模型機率與驗證門檻變化共同影響，因此不將 Recall 的變化單獨解釋為特徵數減少的直接效果。

\section{統計推論範圍與限制}

本章檢定屬於既有結果的事後探索分析，完整保留 165 項比較，避免只報告增加或達顯著者。公司群聚重抽考量同一公司的重複測試紀錄，但使用已固定的模型與驗證決策，不涵蓋重新取得訓練資料、重新選模與其他年度環境的全部不確定性。不同公司也可能受到共同產業或景氣衝擊，這些相依性並未另行建模。

原始 target 為公司最終狀態，且同一公司可能跨擬合與測試期間出現，因此結論限於本資料與既定時間分區下的回溯性公司狀態分類。統計差異不等同模型已具備提前一年預測能力，也不能據此宣稱對全新公司或不同市場具有普遍優越性。
\endgroup
